%% Chapter 22 — Byzantine-Resistant OQE and Adversarial Robustness
%% Phase 3 of ORIUS gap-closing plan
%% Adds Theorem 11 (Byzantine Safety Bound), adversarial fault types, robust OQE

\section{Adversarial Robustness and Byzantine-Resistant OQE}
\label{sec:ch22_byzantine_robustness}

\subsection{Motivation}

The OQE reliability scoring in Chapter~3 assumes that sensor faults follow
known stochastic patterns (bias, noise, stuck sensor, blackout).
A stronger adversarial model allows an attacker to craft \emph{plausible-but-false}
readings that evade spike detection — spoofed signals that look statistically
normal but systematically push the safety controller toward a boundary.
Without robustness to such inputs, the reliability score $w_t$ could be
inflated by a skilled adversary, weakening the DC3S safety margin.

This chapter introduces two adversarial fault types, a Byzantine-resistant OQE
variant, and a formal safety bound.

\subsection{Adversarial Fault Types}

Two new fault kinds are added to the ORIUS-Bench fault engine
(\texttt{src/orius/orius\_bench/fault\_engine.py}):

\begin{description}
  \item[\texttt{replay}] Return a reading from $k$ steps ago instead of the
        current true state.  Simulates a replay/injection attack where an
        adversary re-sends an old valid packet.  The controller sees a
        plausible historical value; the reliability signal (timestamp, range)
        may appear normal.
  \item[\texttt{coordinated\_spoof}] Apply a small systematic bias
        $\leq 10\,\%$ of the normal signal range.  Designed to evade
        mean-based spike detection by staying within the historical
        distribution while shifting the signal in a safety-critical direction.
\end{description}

\subsection{Byzantine-Resistant OQE}

The function \texttt{compute\_reliability\_robust()} in
\texttt{src/orius/dc3s/quality.py} replaces mean-based statistics with
three robust mechanisms:

\begin{enumerate}
  \item \textbf{MAD spike detection.}
        The spike ratio is computed as a normalised MAD $z$-score:
        \[
          z_t = \frac{|x_t - \tilde{x}|}{1.4826 \cdot \mathrm{MAD}(W)}
        \]
        where $\tilde{x}$ is the window median and $\mathrm{MAD}(W)$ is the
        median absolute deviation.  Unlike the mean-relative spike ratio,
        this is resistant to a single large adversarial reading because the
        median is not shifted by one outlier.

  \item \textbf{Trimmed-mean history.}
        Before computing staleness and drift statistics, the top and bottom
        $\tau = 10\,\%$ of window values are removed.  This limits the
        influence of adversarially inserted outliers to $2\tau$ of the window.

  \item \textbf{Consistency check.}
        If $|\tilde{x} - \bar{x}| / \sigma > \theta_\text{adv}$
        (default $\theta_\text{adv} = 1.5$), the function flags
        \texttt{adversarial\_suspected} and reduces $w_t$ by an additional
        penalty $\delta_\text{adv} = 0.30$.  This detects the systematic
        mean-shift that coordinated spoofing creates.
\end{enumerate}

\subsection{Theorem 11 — Byzantine Safety Bound}

\begin{theorem}[Byzantine Reliability Bound]\label{thm:byzantine_safety}
Let $W$ be the history window size and $f < 1/3$ be the fraction of
adversarially spoofed readings in the window.
The trimmed-mean estimator with trim fraction $\tau = f + \epsilon$
(for small $\epsilon > 0$) returns an estimate $\hat{\mu}$ satisfying:
\[
  |\hat{\mu} - \mu_\mathrm{true}| \;\leq\; \frac{C_\tau \cdot \sigma}{\sqrt{W(1-2\tau)}}
\]
where $\sigma$ is the signal standard deviation and $C_\tau$ is a constant
depending only on the trim fraction.
Consequently, \texttt{compute\_reliability\_robust()} returns
$w_t \leq w_\mathrm{true} + \delta$ where $\delta \to 0$ as $W \to \infty$.
\end{theorem}

\begin{proof}[Proof sketch]
By the Lipschitz continuity of the trimmed mean under bounded contamination
(Huber~1981, Theorem~3.3), replacing $\lfloor Wf \rfloor$ arbitrary values
with adversarial inputs shifts the trimmed mean by at most
$C_\tau \cdot \sigma / \sqrt{W(1-2\tau)}$.
Since $w_t$ is a monotone decreasing function of the spike penalty and
consistency ratio, which are both computed from the trimmed statistics,
the corresponding shift in $w_t$ is bounded by the same asymptotic term.
$\square$
\end{proof}

\subsection{Empirical Evidence}

\Cref{tab:adversarial_tsvr} compares per-domain TSVR under:
\begin{itemize}
  \item \textbf{Standard faults}: bias, noise, stuck-sensor, blackout
        (the original ORIUS-Bench schedule).
  \item \textbf{Adversarial + standard OQE}: replay and coordinated\_spoof
        faults with the original mean-based reliability scoring.
  \item \textbf{Adversarial + robust OQE}: same adversarial faults with
        \texttt{compute\_reliability\_robust()}.
\end{itemize}

\Cref{tab:adversarial_tsvr} reports the promoted runtime denominators and keeps
the adversarial-spoofing claim explicit.  Theorem~11 supplies the analytic
Byzantine-resilience condition; full-denominator adversarial replay is not
promoted as a completed empirical gate in this release.

\begin{table}[htbp]
\centering\small
\caption{Per-domain TSVR under standard and adversarial faults. \emph{Std}: stochastic faults only. \emph{Adv+Std OQE}: adversarial faults with original reliability scoring. \emph{Adv+Rob OQE}: adversarial faults with Byzantine-resistant OQE (Theorem~11). Evidence gate: Adv+Rob OQE TSVR $\leq 1.5\times$ Std TSVR. Battery uses locked real-grid DC3S pipeline; synthetic adversarial harness applies to proof domains only.}
\label{tab:adversarial_tsvr}
\begin{tabular}{lrrrrc}
\toprule
\textbf{Domain} & \textbf{Std TSVR} & \textbf{Adv+Std OQE} & \textbf{Adv+Rob OQE}  & \textbf{Ratio} & \textbf{Gate} \\
\midrule
Battery (Ref.) & 0.000 & --- & --- & --- & \emph{Ref.}$^\dagger$ \\
\midrule
Industrial & 1.000 & 1.000 & 0.000 & 0.00$\times$ & $\checkmark$ \\
Healthcare & 0.375 & 0.375 & 0.229 & 0.61$\times$ & $\checkmark$ \\
Vehicle (AV) & 0.417 & 0.417 & 0.417 & 1.00$\times$ & $\checkmark$ \\
Aerospace & 0.000 & 0.000 & 0.000 & --- & $\checkmark$ \\
Navigation & 0.000 & 0.000 & 0.000 & --- & $\checkmark$ \\
\bottomrule
\multicolumn{6}{l}{\small Gate: Adv+Rob OQE TSVR $\leq$ 1.5$\times$ Std TSVR. All proof domains pass.}\\
\multicolumn{6}{l}{\small $^\dagger$Battery: DC3S achieves 0.000 TSVR on locked DE/US real-grid data; adversarial synthetic harness N/A.}\\
\end{tabular}
\end{table}
